\subsection{Four-transvection regime theorem}
\label{sec:bt1249-four-transvection-regimes}

BT1236 establishes the exact four-transvection Clifford target and BT1233 word-metric fingerprint.  BT1242--BT1248 sharpen this into a regime theorem over all four-projective-transvection sets.

There are
\[
\binom{40}{4}=91390
\]
four-projective-transvection sets.  Exactly
\[
61560
\]
of them generate the full group \(Sp(4,3)\).  Full-order generation therefore is not rare, but it is not sufficient as a tomography invariant.

The full-order sets split into three word-metric diameter regimes:
\[
51840_{\operatorname{diam}=10}^{22680},
\qquad
51840_{\operatorname{diam}=12}^{25920},
\qquad
51840_{\operatorname{diam}=14}^{12960}.
\]
The BT1228/BT1233 fingerprint is the diameter-14 regime, hence it occupies
\[
\frac{12960}{61560}=0.2105263158
\]
of the full-order four-transvection sets.

The local closure laws distinguish these regimes.  The diameter-10 regime is the fast regime in which all four triples already close to order \(648\).  The diameter-12 regime has pair pattern
\[
9^2\,24^4
\]
and triple pattern
\[
72^1\,648^3.
\]
The BT1228/BT1233 diameter-14 regime has the balanced pair pattern
\[
9^3\,24^3
\]
and balanced triple pattern
\[
72^2\,648^2.
\]
Thus the diameter-14 fingerprint is not selected by maximal stabilizer size.  It is selected by balanced local closure: neither every triple is already large, nor does the set fail to reach the full group.

\paragraph{Tomography consequence.}
A recovered four-transvection model must pass three increasingly strict gates: closure order \(51840\), correct diameter \(14\), and the BT1233 sphere/ball-growth fingerprint.  The BT1242 split proves that the first gate alone admits wrong full-order recoveries.
